\chapter{Corrections, retractions and withdrawn readings}\label{ch:corrections}

This appendix collects the corrections, retractions, errata, withdrawn
claims and reversed verdicts recorded in the program's evidence directories
and ledgers. Where a chapter tells a correction at length, the entry here
is a one-line pointer. Chapter~\ref{ch:audit}
and \S\ref{sec:corrections} hold the corrections that moved index calculus
against rho on the leaderboard (the retired online wins, the priced word
XORs, rule~v8, the cross-host F4 counter drift, the volcano-crater and
plane-quartic withdrawals); the rho-side corrections are told in
Chapters~\ref{ch:matchedrho} and~\ref{ch:rhovariants}, and the
unconventional-route corrections in Chapters~\ref{ch:tensor}
and~\ref{ch:learning}. They are not repeated here.

Each entry gives what was first read, the record that read it, what replaced
the reading and when, and the source. Three conventions apply throughout.
\begin{itemize}
\item \emph{Standing.} Unless a row says \emph{validated}, the correction is
  itself a producer's record, dated and additive (the frozen evidence it
  corrects is left in place), and carries no proof-carrying-claim verdict.
  Under the rule that every producer may be sabotaging the work, such a row
  is an unverified assertion of its producer, as is the reading it replaced.
\item \emph{Direction.} A correction of the rho reference or of the pricing is
  an accounting change. None of the rows below is an improvement of index
  calculus, and none moves a validated leaderboard row below its reference.
\item \emph{Cross-references.} Seven corrections are carried in the main
  text (Table~\ref{tab:corr-xref}) and listed here with a pointer only.
\end{itemize}

\begin{table}[tbp]
\centering\small
\begin{tabularx}{\linewidth}{@{}>{\raggedright\arraybackslash}p{0.2\linewidth}>{\raggedright\arraybackslash}X>{\raggedright\arraybackslash}X>{\raggedright\arraybackslash}p{0.17\linewidth}@{}}
\toprule
Item & Original reading & Correction & Main text \\
\midrule
H-PSEUDO character-sum constant & Fourier-flatness offset $R\approx1.28$--$1.34$, $C(p)\sim p^{0.055}$--$p^{0.079}$, open & 4 Nov: the $(x,y)$-keyed small-$x$ base is negation-closed and the null was not; against a negation-closed null $R$ is consistent with 1 (pooled 56 curves $1.005\pm0.008$); verdict \emph{explained} under a non-rejection rule, so a residual up to about $+9\%$ at $2^{20}$ is not excluded & Chapter~\ref{ch:summer} \\
\addlinespace
Xedni lifting & ``closed'' in exact form & 1 Nov: the Jacobson--Koblitz--Silverman--Stein--Teske obstruction is conditional (Lean audit); overbroad readings refuted; algorithm-specific bridges still open & Chapters~\ref{ch:obstructions}, \ref{ch:summer} \\
\addlinespace
$\ell=8$ hybrid-guess cells & the $\ell=8$ guess run had not completed & 22 Oct: four guessed variables bring every $\ell=8$ draw back to degree 4; 23 Oct: three do not, for that guess set & \S\ref{sec:binary-descent}, \S\ref{sec:frontier} \\
\addlinespace
Planner weight rules & \texttt{ic-planner/2}, \texttt{/3}: $1.41N_0$ (prime, $w=1$), $0.71c_0$ (Koblitz, $w\ge0.3$) & 26 Oct: both failed their preregistered out-of-sample tests; \texttt{ic-planner/4} withdraws them & \S\ref{sec:prime} \\
\addlinespace
EcGFp5 symmetrisation & no $m=5$ ideal degree measured & 1 Nov: $D=2^{16}=65{,}536$ at all three primes & Chapter~\ref{ch:summer} \\
\addlinespace
$n=83$ Weil-descent refutations & F4, F5, F6-IC, crossbred, SAT-descent and MQ-FES ``proved UNSAT'' on planted relations & 24 Oct: \texttt{FieldStructure} truncated the field table to one 64-bit word; builders now refuse such fields; every earlier $n=83$ engine refutation is void & \S\ref{sec:bo-n83} (summary in \S\ref{sec:n83}) \\
\addlinespace
Joux--Vitse sieve arm & 19 Oct comparisons withheld by the 20 Oct review; replication failed its rho A/A rule & 2--3 Nov: matched negation-map rho at $p=1511$; sieve below rho in every cell of the band; \emph{validated} (\texttt{VAL-20261103-c31c76}); no pcc claim, so its statement is an unverified assertion & \S\ref{sec:transfer} \\
\bottomrule
\end{tabularx}
\caption[]{Corrections carried into the main text, listed here by pointer.
Sources: \art{crypto-autoresearcher:docs/reports/hpseudo-symnull-20261104/README.md};
\art{research/jksst_xedni_audit_20261101/};
\art{research/dreg_closure_guess_20261022/RESULTS.md},
\art{research/dreg_closure_guess3_20261023/RESULTS.md};
\art{research/ic_planner_weight_rule_20261021/README.md},
\art{research/ic_planner_weight_rule_binkob_20261026/README.md};
\art{research/ecgfp5_quotient_solve_20261101/README.md};
\art{research/n83_msweep_20261024/README.md};
\art{research/jv_sieve_negrho_20261102/}.}
\label{tab:corr-xref}
\end{table}

\section{Benchmark rows removed from the admitted table}\label{sec:corr-archive}

The validation gate (19 October) froze 116 benchmark rows as unvalidated
(pre-gate). By 6 November 65 rows had been removed from the admitted table
to \art{benchmarks/ARCHIVE.json}, each with its
exact former text, its SHA-256 and a disposition (Table~\ref{tab:fm-archive}).
A row is \emph{incomplete} when its evidence lacks a replay certificate, a
validator report or a matched rho, \emph{documentary} when it records a
paper, census or synthesis rather than a measurement, and \emph{censored} when it rests on
censored sessions and lower-bound costs. All 65 rows date from 13--26 October; 47 of them
belong to the $n=83$ pipelines, the isogeny-volcano studies and the IC tool
programme's exploration rounds. Removal is not a refutation: the archived
text is kept as historical evidence of what was claimed.

\begin{table}[tbp]
\centering\footnotesize
% generated by make_tab_fm.py; do not edit
\begin{tabular}{@{}lrrrr@{}}
\toprule
entries & incomplete & documentary & censored & total \\
\midrule
$n=83$ pipelines & 20 & 2 & 0 & 22 \\
isogeny volcanoes and classes & 10 & 0 & 0 & 10 \\
IC tool programme rounds (A14--A20, R06--R12, scan) & 13 & 0 & 0 & 13 \\
other & 13 & 6 & 1 & 20 \\
\midrule
all & 56 & 8 & 1 & 65 \\
\midrule
entry date & incomplete & documentary & censored & total \\
\midrule
2026-10-13 & 0 & 2 & 0 & 2 \\
2026-10-14 & 0 & 2 & 0 & 2 \\
2026-10-15 & 1 & 1 & 0 & 2 \\
2026-10-16 & 5 & 0 & 0 & 5 \\
2026-10-17 & 22 & 3 & 0 & 25 \\
2026-10-18 & 16 & 0 & 0 & 16 \\
2026-10-19 & 8 & 0 & 1 & 9 \\
2026-10-20 & 1 & 0 & 0 & 1 \\
2026-10-21 & 1 & 0 & 0 & 1 \\
2026-10-26 & 2 & 0 & 0 & 2 \\
\bottomrule
\end{tabular}

\caption[]{Benchmark rows removed from the admitted table, by topic and by
the date in their entry name. Topics are assigned by directory-name
keyword in \art{figures/make_tab_fm.py} (first match wins); counts from
\art{benchmarks/ARCHIVE.json} via \art{figures/data-fm-glance.json}.}
\label{tab:fm-archive}
\end{table}

\section{Binary and Koblitz fields}\label{sec:corr-binary}

Table~\ref{tab:corr-binary} lists the binary-field items. Two need a
sentence beyond the table.

\paragraph{Torsion anchors.} The merged report of
\art{research/ic_m2_cycle_leads_20261030} (PR~\#3009) stated that torsion
anchors ``cost $1.5\times$ ($n=23$) to $16\times$ ($n=31$) on cofactor-4
curves''. A frozen-protocol test at $n=41$ found the on/off trial ratio
$0.975$ $[0.641,1.420]$, and a post-hoc replicate of 32 paired targets at
$n=31$ gave $13.3\times$ $[10.3,17.5]$: the penalty exists at $n=31$, is about
$1.5\times$ at $n=23$ and is absent at $n=41$, so the range is not to be
quoted as a property of torsion anchors. A companion study on $a=0$, $n=37$
found a penalty of $6.3$--$14.5\times$ at every cofactor-to-base ratio from
$0.30$ to $10.84$. Its README also records that the earlier bootstrap
intervals (PRs \#3126 and \#3138) were drawn with \texttt{mawk}'s
non-reproducible \texttt{srand}; recomputed with a seeded generator they move
by at most $0.02$. The earlier $n=37$ figure $0.134$ $[0.082,0.205]$ was
off/on; it is the same measurement as the on/off $7.46$ $[4.90,11.92]$ and is
not a reversal
(\art{research/ic_m2_torsion_n41_20261101/REPORT.md},
\art{research/ic_m2_torsion_ratio_n37_20261101/REPORT.md}).

\paragraph{The $k=2$ collision excess.} On $y^2+xy=x^3+1$ over
$\FF_{2^{19}}$ ($r=130{,}873$, birthday expectation $5.004$) the archived
census read informative-to-birthday ratios of $3.397$ and $7.593$ at $k=2$
and called them a $k$-specific collision excess. The census counted a
colliding pair as informative unless its difference lay in the span of one
kernel vector per generator, but the factor base is the union of
$\langle\tau,-1\rangle$ orbits, so negation identities and $\tau$-shifted
kernel relations were counted as informative. Replaying the archived code
classifies 16 of 17 and 33 of 38 pairs as structural; the genuine-to-expected
ratios are then $0.20$ and $1.00$ at $k=2$ and $0.40$ and $0.80$ at $k=3$, none
above the band
(\art{crypto-autoresearcher:analysis/ecdlp2m-iteration-20261009/certbin-k2-collision-audit/README.md}).

\begin{table}[tbp]
\centering\small
\begin{tabularx}{\linewidth}{@{}>{\raggedright\arraybackslash}p{0.17\linewidth}>{\raggedright\arraybackslash}X>{\raggedright\arraybackslash}X>{\raggedright\arraybackslash}p{0.21\linewidth}@{}}
\toprule
Item & Original reading & Correction & Source \\
\midrule
Torsion anchors, cofactor 4 & PR~\#3009: $1.5$--$16\times$ more trials & 1 Nov: $0.975$ $[0.641,1.420]$ at $n=41$; not a general property (see text) & \art{research/ic_m2_torsion_n41_20261101/} \\
\addlinespace
Semaev 2015, $t=2$ & first fall degree ``always 4'' (\S4.5); corpus inference that these rows refute $d_{F4}\le d_{ff}$ & 2 Nov: $d_{ff}=2$ on 200 of 200 instances by a trace certificate; the solver needs degree 2 on 154 and 3 on 46, never 4, so the sharp assumption fails on 46/200 only & \art{research/semaev_t2_trace_20261102/REPORT.md}; cf.\ Corollary~\ref{cor:assump} \\
\addlinespace
Hamming ball $B_3$, $n=19$ & $R=1.10$ against random subspaces & $1.10$ came from one noisy subspace denominator; against 8 matched subspaces $R=0.27$ $[0.13,0.46]$ and $0.32$ $[0.14,0.55]$; \emph{not decided} & \art{crypto-autoresearcher:knowledge/leads/ecdlp/ic-binary/LEAD-ICB-nonlinear-frob-stable-bases.yaml} \\
\addlinespace
SAT answers in \texttt{EV-ICPERF-390707} & five unverified SAT answers & they are decompositions on the quadratic twist $y^2+xy=x^3+1$; instrument not invalidated, SAT-over-Gr\"obner gap direction reproduced & \art{crypto-autoresearcher:ledger/evidence/EV-ICPERF-390707.yaml} \\
\addlinespace
Certified $k=2$ collisions & 17 and 38 pairs against $5.004$: a $k$-specific excess & instrument artefact; 16/17 and 33/38 structural (see text) & certbin audit (see text) \\
\addlinespace
Learned factor-base selection & producer: ``about 9 bits, identical across the band'' behind rho & 5--6 Nov: a pinned-calibration point; ridge pick $1.8/14.8/17.7$ bits ($n=17$) and $3.3/15.1/17.9$ ($n=19$) behind strong rho at $w=0/0.1/1$; the lead records independent validation & \art{crypto-autoresearcher:knowledge/leads/ecdlp/ic-binary/LEAD-ICB-learned-cost-model-factor-base-selection.yaml} \\
\addlinespace
Root linear-row harvest (\#1862) & raised counted word XORs about 7\% (one workload, $K_0$, $n=13$) & 24 Oct: the sign depends on size: $S$ up $9.6\%$ and $7.0\%$ at $n=11,13$, down $0.3\%$ and $4.1\%$ at $n=17,19$; every priced F4 arm stays $7$--$1{,}250\times$ strong rho & \art{research/f4_harvest_remeasure_20261024/README.md} \\
\bottomrule
\end{tabularx}
\caption[]{Binary and Koblitz corrections not in the main text. All are
producer-reported except the learned-selection row, whose atlas lead records
a \texttt{VALIDATION.json}; the treatise did not open that repository. The
already-recorded binary corrections (September reversals, the retired $n=61$
$13.94\times$ and $n=83$ $18.53\times$ online readings, Table~A's understated
gap, the contradicted degree-4 floor) are in Chapter~\ref{ch:audit} and
\S\ref{sec:corrections}.}
\label{tab:corr-binary}
\end{table}

\section{Prime fields}\label{sec:corr-prime}

Table~\ref{tab:corr-prime} lists the prime-field items; the first needs a
paragraph.

\paragraph{Byte counters priced as operations.} The \S5a band priced
\texttt{residual\_table\_\allowbreak bytes\_\allowbreak uncharged}, a byte count of the residual
table, at $w$ GAE per byte. On the residual arms it was $83$--$99\%$ of $S$ at
$w=1$, and it alone made their $w>0$ distance to rho $17$--$58\times$ rather
than $2.6$--$5\times$. On 1~November the owner chose option~B: the band prices
operations only, and \texttt{band.rs} lists the byte counters without pricing
them. Re-running every committed weighted comparison under the rule turned all
24 residual-versus-planner comparisons of \art{research/ic_prime_residual_lp_20261026/}
from ``unresolved inside the band'' to ``residual cheaper at every weight''. It
also resolved 36 of 136 comparisons in the planner-exhaustive study and 2 of 6
in the Koblitz weight-rule study, and reversed no direction. Its effect cut both ways: the
pre-registered criterion C1 of \art{research/ic_residual_scale_20261101/} is met
in 2 of 8 cells, not 6 of 8, because the \texttt{eq} candidate's apparent gain
at $w=1$ rested on byte counts. The $w=0$ results and the $2.59\times$
leaderboard row are unchanged
(\art{research/ic_byte_counter_pricing_20261101/README.md}).

\begin{table}[tbp]
\centering\small
\begin{tabularx}{\linewidth}{@{}>{\raggedright\arraybackslash}p{0.17\linewidth}>{\raggedright\arraybackslash}X>{\raggedright\arraybackslash}X>{\raggedright\arraybackslash}p{0.21\linewidth}@{}}
\toprule
Item & Original reading & Correction & Source \\
\midrule
Residual collectors at $w>0$ & res4/rho $41$--$58\times$ at $w=1$; C1 met in 6 of 8 cells & 1 Nov, option~B: res4/rho $3.6$--$4.8\times$; C1 met in 2 of 8 & \art{research/ic_byte_counter_pricing_20261101/} \\
\addlinespace
Pooled $k$-sum tables & no \texttt{VALIDATION.json}; unqualified & 12 Nov: gate-validated (16 rows), claim \texttt{PCC-20261112-31adb2}, unverified standing; \S\ref{sec:pg-ksum} & \art{research/prime_ksum_pooled_20261029/VALIDATION.json} \\
\addlinespace
Multi-target IC (ic4) & cheaper than rho at $T=16384$ (\texttt{PCC-20261111-fa13c8}) & 11 Nov: withdrawn as a statement about rho after the rho start fix; superseded by \texttt{PCC-20261111-9a26ac}; \S\ref{sec:pg-multitarget} & \art{research/mt_rho_start_fix_20261111/} \\
\addlinespace
Tensor-contraction route & ``nothing was measured'' & 11--12 Nov: graph widths measured; exact Boolean contraction of the encodings costs at least about $2^{2b}$; no cost of the route measured; \S\ref{sec:pt-tensor} & \art{research/semaev_tensor_treewidth_20261111/} \\
\addlinespace
Small-root lattice reach & $\varepsilon(\Sem3)=0.8408$, best margin $-0.0481$ at $m=4$ & 24 Oct (\texttt{CORR-20261024-b2a799}): wrong exponent in \texttt{smallroot\_reach.py}, $\varepsilon_{HG}=\varepsilon_{TOOL}/n$; $\varepsilon(\Sem3)=0.2803$, margin $-0.6086$ at $(4,2)$, best $-0.5161$ at $m=32$ & \art{crypto-autoresearcher:analysis/smallroot-reach/CORRECTION.md} \\
\addlinespace
Oracle gate of the Semaev-alternatives memo & ``$e<1$ at $m=4$'' & 4 Nov: not implied by its own cost model; a generic 2+2 MITM passes it ($s=1.93$) at cost $p^{2/3}$; correct gate $s<m-3$, i.e.\ $e<(m-3)/4$ & \art{research/semaev_alternatives_audit_20261104/README.md} \\
\addlinespace
Cofactor-cleared prime base & $G_0=6.180$ $[5.915,6.456]$, pass line met at every weight & 3 Nov: verdict \texttt{invalid\_controls} by the frozen rule (positive control's lower limit $1.149<1.2$); not to be cited as a gain & \art{research/ic_cofactor_clear_ab_20261103/README.md} \\
\addlinespace
Multi-target Koblitz figures & closest batch figure $1.17\times$ rho & accounting: $6.51\times$ batch rho at $k=32$ (Kuhn--Struik reference), $6.38\times$ with step and build priced; one target $43.7\times$ & \art{research/ic_rho_koblitz_20260923/README.md} \\
\bottomrule
\end{tabularx}
\caption[]{Prime-field corrections, with one Koblitz accounting row that
belongs with them because it changed a rho reference rule. All are
producer-reported; the small-root correction is a single-agent computation
that has had no independent review.}
\label{tab:corr-prime}
\end{table}

\section{Engines and measurement}\label{sec:corr-engine}

Several corrections changed the instrument, not a result
(Table~\ref{tab:corr-engine}). Five are told at length in
Chapter~\ref{ch:matchedrho} and only named here: calibration round 1
(\S\ref{sec:mr-calibration}), the independent-validation host class and
the statistics~v2 intervals with censored rho (\S\ref{sec:mr-stats}), the
bounds frontier v1$\to$v2 (\S\ref{sec:mr-frontier}), and the $n=29$
generator-orbit panel, which is a Frobenius-folded BSGS
(\S\ref{sec:gf-bsgs}). They matter
because intervals and references quoted earlier in this treatise were
computed under the older rules.
\begin{itemize}
\item \emph{Accounting v2} (23 October) sized the strong Koblitz reference
  to $r$. Its mean $S$ falls to $0.182$ $[0.171,0.196]$, $0.223$
  $[0.201,0.248]$ and $0.678$ $[0.499,0.891]$ of the frozen id at $n=13$, 19
  and 37, and its fixed set-up share from $83.3\%$, $71.3\%$ and $30.2\%$ to
  $39.9\%$, $29.3\%$ and $4.3\%$. Any small-$n$ comparison against the frozen
  strong id read rho too expensive by up to that factor. The report states
  that this is a correction of the reference instrument, not a faster
  algorithm (\art{research/ecbench_accounting_v2_20261023/REPORT.typ}).
\item \emph{In-house F4 is not a published engine.} A desk review found that
  the repository's F4 needs $10^3$--$10^7$ more operations than a brute grid
  on the same PKM systems, so root-finding-over-F4 ratios such as those of
  \S\ref{sec:rootfinding} (for example $383\times$) are not evidence about
  published Gr\"obner solvers; degree measurements made with msolve are
  unaffected (\art{research/ic_sweeper_20261024/README.md}).
\end{itemize}

\begin{table}[tbp]
\centering\small
\begin{tabularx}{\linewidth}{@{}>{\raggedright\arraybackslash}p{0.17\linewidth}>{\raggedright\arraybackslash}X>{\raggedright\arraybackslash}X>{\raggedright\arraybackslash}p{0.21\linewidth}@{}}
\toprule
Item & Original reading & Correction & Source \\
\midrule
Strong Koblitz rho id & frozen id, set-up $30$--$83\%$ of $S$ & 23 Oct: v2 sized to $r$, $S$ ratio $0.182$--$0.678$ (see text) & \art{research/ecbench_accounting_v2_20261023/} \\
\addlinespace
Runner PR \#2455 & IC online window ``about 15\% lower'' & 23 Oct: counts identical in 10 of 10 pairs; online ratio $1.067$, not distinguished; the $18.4$--$38.8\times$ shorter small-arm walls are an allocator effect (\texttt{malloc\_trim}, \texttt{/proc} reads) & \art{research/ecbench_allocator_confound_20261023/} \\
\addlinespace
E42 Macaulay tails & preregistered control range $[900,960]$ re-allocations & 3 Nov: the range was mis-derived (measured $245.6$); E42's kill fired with $W$ unranked; E42b re-registered before any ranking data & \art{research/e42_macaulay_tails_20261103/REPORT.typ} \\
\bottomrule
\end{tabularx}
\caption[]{Instrument and measurement corrections not told elsewhere. All
producer-reported; none changes an IC count except where the row says so.
Calibration round 1, independent validation, statistics v2 and the bounds
frontier v1 are in Chapter~\ref{ch:matchedrho}; the $n=29$ panel is in
\S\ref{sec:gf-bsgs}.}
\label{tab:corr-engine}
\end{table}

\section{Rho references}\label{sec:corr-rho}

Table~\ref{tab:audit-rho} gives the rho-side changes listed with the
leaderboard corrections. Table~\ref{tab:corr-rho} adds those that moved
study readings and are not told elsewhere; the negation-map audit
(\S\ref{sec:mr-negaudit}), the withdrawn Cheon/rho readings on the P-224
shape (\S\ref{sec:gf-cheon}), the Curve25519 \texttt{wwalk} re-baseline
(\S\ref{sec:gf-collision}) and the hyperelliptic rho reference
(\S\ref{sec:rv-jacobian}) are in
Chapters~\ref{ch:genericfloor}--\ref{ch:rhovariants}. One
binary defect underlies the first two rows: the build that wrote the
crater-automorphism and prime-automorphism sessions left the
\texttt{rho.cubic}/\texttt{rho.quartic} class set-up ($16$--$58$ GAE per
solve) out of each record's total
(\art{research/class_setup_impact_20261024/}).

\begin{table}[tbp]
\centering\small
\begin{tabularx}{\linewidth}{@{}>{\raggedright\arraybackslash}p{0.17\linewidth}>{\raggedright\arraybackslash}X>{\raggedright\arraybackslash}X>{\raggedright\arraybackslash}p{0.21\linewidth}@{}}
\toprule
Item & Original reading & Correction & Source \\
\midrule
Cheon with auxiliary inputs & no exponent separation established; medium-tier re-run frozen, no results & 7--11 Nov: the validated 11-size study separates the fitted exponents over 36--56 bits in both families ($r-1$ by $0.002$, $r+1$ by $0.026$--$0.032$); the medium tier reproduced 272 of 272 runs; \S\ref{sec:gf-cheon} & \art{research/cheon_ecbench_wide_20261106/REPORT.md} \\
\addlinespace
Automorphism rho, prime & cubic/negation $0.611$, quartic/negation $0.791$ & 24 Oct: $0.647$ $[0.543,0.770]$ and $0.819$ $[0.722,0.929]$; the quartic interval no longer contains $1/\sqrt2$, so finding 1 holds for cubic only & \art{research/prime_rho_automorphism_20261015/README.md} \\
\addlinespace
Volcano crater controls & $j=0$ $0.697$, $j=1728$ $0.707$ & 24 Oct: $0.716$ $[0.598,0.857]$ and $0.723$ $[0.621,0.843]$; both still exclude 1 & \art{research/isogeny-volcano-cost_20261017/README.md} \\
\addlinespace
GLV-invariant matched walk (E27) & distances of \S9.24 & 30 Oct: every row $1.448$--$1.456\times$ further from rho at $w=0$; \emph{validated}; told in \S\ref{sec:rv-aut} & \art{research/e27_walk_20261030/REPORT.typ} \\
\addlinespace
Order-24 walk (E29) & crossing with order 6 at $w=1$ (\S9.33) & with its $40.6$ unpriced mod-$n$ operations per step priced, order 24 loses at every weight; the crossing was a one-op-per-step artefact & GLV note, \S9.39 (E34) \\
\addlinespace
Verschiebung split & round 1: $\lambda=t-1$ via $V^n$, ratio $0.49$, ``not established'' & round 2: $V^n=[2^n]\circ\pi^{-1}$ is $n$ doublings plus square roots; ratio $\ge2.5$ at $c_{\rm dbl}=8$, a provable loss & \art{research/isogeny_loop_glv_20261011/README.md} \\
\bottomrule
\end{tabularx}
\caption[]{Corrections to rho references used by studies. GLV note:
\art{research/notes/index-calculus/RESEARCH_GLV_INVARIANT_FACTOR_BASES.md}, whose
\S9.19 also records seven textual errata found by the validation gate with no
result, verdict or ranking changed. All rows except E27 are producer-reported.}
\label{tab:corr-rho}
\end{table}

\section{Transfers and covers}\label{sec:corr-transfer}

The transfer audit's v2 certificates (amendment 3, 26 October) changed the
baseline itself. The model's $\sqrt{\pi r/4}$ is already rho with negation;
dividing it again by $\sqrt c$ counted negation twice and put every baseline
$0.5$ bit too low. The amendment also gave rho on a binary cover's Jacobian
image the Frobenius class of the baseline curve (Koblitz cover ratios had
been overstated by $\tfrac12\log_2(n/d)$ bits), required a margin of 8 bits
before a price counts as below the baseline, and kept the verdict
\texttt{refuted} only for a carried route priced as a route; formula sweeps,
identity-path prices and unconstructed maps are now \texttt{unknown}. The
regeneration first overwrote the v1 certificates; they were restored byte for
byte in \texttt{v1/} on 27 October
(\art{research/transfer_audit_20261011/PROTOCOL.md}). The other transfer
items are in Table~\ref{tab:corr-transfer}.

\begin{table}[tbp]
\centering\small
\begin{tabularx}{\linewidth}{@{}>{\raggedright\arraybackslash}p{0.17\linewidth}>{\raggedright\arraybackslash}X>{\raggedright\arraybackslash}X>{\raggedright\arraybackslash}p{0.21\linewidth}@{}}
\toprule
Item & Original reading & Correction & Source \\
\midrule
Transfer-audit baseline & \S4 baseline; \texttt{refuted} for any route not below it & 26 Oct: every baseline $+0.5$ bit; 8-bit margin; \texttt{refuted} narrowed (see text) & \art{research/transfer_audit_20261011/} \\
\addlinespace
Cover descent at the genus floor & 846 \texttt{cover\_descent} rows \texttt{formula\_advantage}; earlier reading: all rest on a genus no GHS cover has & priced at the genus the construction has: 599 remain (576 at $k=6$, 23 at $k=14$), 41 marginal, 206 boundary; rows stay \texttt{proposed} & \art{research/cover_descent_genus_floor_20261015/README.md} \\
\addlinespace
Hilbert stop on the weak-class cover (B0) & counted $S$ falls $2.53$--$2.64\times$ & 23 Oct: relabelling; unpriced F4-driver operations rose $8.3$--$8.5\times$; the repaired form B is engineering ($1.85$--$1.88\times$) & \art{research/ic_cover_weakclass_20261023/REPORT.md} \\
\addlinespace
GHS cover census, $n=4$ & preregistered: floor 21 for $E(K)[2]=\ZZ/2$; odd-class minima in $\{13,25\}$ & both failed; derived floor 25; odd class 13 iff $j\notin\FF_{q^2}$, else 37, never 25; one check bug fixed & \art{research/ghs_cover_census_n4_20261024/PROTOCOL.md} \\
\addlinespace
Isogeny seed, curve25519 & seed 1093 cheaper than seed 0 & 3 Nov: V2 failed (mean $S$ ratio $0.641>0.6$); interval $[0.379,1.084]$ includes 1; gain not established & \art{research/curve25519_isogeny_seed_20261103/README.md} \\
\addlinespace
CM canonical lift (lane L9) & closure argument as first written & \texttt{CORR-20260810-930c42}: closes only ordinary prime-field Serre--Tate/global CM models of fixed arity $2$--$9$; point-lifting routes give $\Pr\le p^{-1+o(1)}$, exponent 1 against rho's $1/2$ & \art{crypto-autoresearcher:knowledge/leads/ecdlp/transfer/LEAD-TRF-cm-canonical-lift-heights.yaml} \\
\addlinespace
GHS magic number & legacy \texttt{ec\_trapdoor::\allowbreak magic\_number} (plain orbit) & 13 Oct: MMT definition; consequences for ECC2K-130 are in \S\ref{sec:koblitz} and \S\ref{sec:ladder3} & \art{research/weak_isogeny_classes_20261013/REPORT.md} \\
\addlinespace
JV sieve, table-line arm U & ``no validation file'' (\S\ref{sec:jac-jv}, through 11 Nov) & validated 11 Nov, 02:01Z: U at $0.0141$, $0.0857$, $0.0900$ of matched rho at $w=0$, $0.1$, $1$ (weak class only); text corrected 12 Nov & \art{research/jv_sieve_line_vs_rho_20261103/VALIDATION.json} \\
\addlinespace
JV Nagao arm against matched rho & $1.41$ and $1.16$ of matched rho; $1.59$--$1.98$ against a pooled constant & 12 Nov: with the Hilbert stop, $1.003$ $[0.916,1.104]$, $0.837$, $0.695$ at $p=1301$, 1451, 1601 (local, unverified); \S\ref{sec:ext-jv-nov} & \art{research/jv_nagao_hilbert_20261112/rescore/table.md} \\
\bottomrule
\end{tabularx}
\caption[]{Corrections in the transfer and cover studies. All
producer-reported.}
\label{tab:corr-transfer}
\end{table}

\section{Structure: conductor gaps, curve traits and instruments}\label{sec:corr-structure}

\paragraph{Conductor gaps.} The conductor-gap series has the longest
correction trail. The CM-pairs study first stated that ``the Lean-provable
theorem'' gave the null for generic methods. On 16 October a sheriff review
(\texttt{SHF-20261016-3caa6f}, finding L5) noted that
\texttt{lean/GenericInvariance.lean} covers only algorithms whose control
flow depends on equality answers and a random tape, that is BSGS-type
algorithms, while ecbench's rho and kangaroo hash the point representation.
On 17 October a second review (\texttt{SHF-20261017-e4d7f6}, L-N2) withdrew
the conclusion ``no implemented method sees the gap'', and on 28 October
(\texttt{SHF-20261028-0f3f64}, PPN-1) the banner that still carried it was
corrected. The open rho reading at $\ell=8209$, $1.22$ $[1.005,1.46]$ from 4
targets, was then excluded by 12{,}288 paired solves under a fresh seed: any
floor/crater difference is below about $1.2\%$ there. In the easy-gap tier
($\ell\approx2^{25}$) the frozen interval $0.917$ $[0.851,0.988]$ excluded 1,
but a preregistered twin round put the decisive floor-twin/crater-twin ratio
at $1.040$ $[0.966,1.119]$: the shift attaches to the coordinate models, not
to the endomorphism ring
(\art{research/conductor_gap_cm_pairs_20261010/README.md},
\art{research/conductor_gap_rho_large_sample_20261018/README.md}). The
index-calculus side is in \S\ref{sec:transfer}; the remaining structural
items are in Table~\ref{tab:corr-structure}.

\begin{table}[tbp]
\centering\small
\begin{tabularx}{\linewidth}{@{}>{\raggedright\arraybackslash}p{0.17\linewidth}>{\raggedright\arraybackslash}X>{\raggedright\arraybackslash}X>{\raggedright\arraybackslash}p{0.21\linewidth}@{}}
\toprule
Item & Original reading & Correction & Source \\
\midrule
Conductor gap, CM pairs & ``no implemented method sees the gap''; null is a theorem for generic methods & 16, 17, 28 Oct: theorem covers BSGS-type arms only; banner withdrawn and qualified (see text) & \art{research/conductor_gap_cm_pairs_20261010/} \\
\addlinespace
Rho at $\ell=8209$; easy-gap tier & $1.22$ $[1.005,1.46]$; post hoc $1.17$ $[1.00,1.38]$ & 18 Oct: first excluded; second attributed to coordinate models ($1.040$ $[0.966,1.119]$) & \art{research/conductor_gap_rho_large_sample_20261018/} \\
\addlinespace
Round-1 crater/floor set-up gap & prime floor/crater $S$ ratio $3.10$ & round 2, A0: a closure counting defect (the floor paid $3.56\times$ the needed set-up); ratio $1.00$ & \art{research/isogeny_conductor_gap_round2_20261011/README.md} \\
\addlinespace
Drastic prime gap, $j=0$ & floor/crater rho quotient $0.5286$ & round 2, 32 fresh targets: $0.9468$ $[0.6622,1.3474]$ (99\%); target-sampling noise & \art{research/drastic_prime_conductor_gap_round2_20261012/README.md} \\
\addlinespace
NIST skip criterion & coefficient test separates the skipped curves & 25 Oct: the smaller P-224 root rejects the published model, and the two models are isomorphic; representation-dependent, not an intrinsic weakness & \art{research/nist_skip_criterion_20261025/REPORT.md} \\
\addlinespace
Coordinate census \texttt{EV-ECDLP-81f4d4} & strength \emph{strong} & 26 Oct (\texttt{CORR-20261026-c94668}): lowered to \emph{preliminary}, with other records of the same sheriff case & \art{crypto-autoresearcher:ledger/corrections/CORR-20261026-c94668.yaml} \\
\addlinespace
Weak-class instrument I1 & \emph{invalid} on all three classes & 4 Nov: the A/A gate (5\% on 20+20 draws) was miscalibrated; with a calibrated gate I1 is valid and blind; the paired-curve anomaly did not replicate (108+108 curves) & \art{crypto-autoresearcher:docs/reports/weakclass-i1-20261104/}, \art{crypto-autoresearcher:docs/reports/paired-anomaly-20261104/} \\
\addlinespace
Polynomial-ring form (FHJRV) & ``unavailable'' unless $b$ is a square & 2 Nov (\texttt{CORR-20261102-7e3a91}): holds for the $x$ (Kummer) class only; for $u=x/y$ on a double-odd curve the quotient is a polynomial ring & \art{crypto-autoresearcher:ledger/corrections/CORR-20261102-7e3a91.yaml} \\
\bottomrule
\end{tabularx}
\caption[]{Structural corrections. All producer-reported, the sheriff
findings included.}
\label{tab:corr-structure}
\end{table}

\section{Unconventional routes}\label{sec:corr-exotic}

\begin{table}[tbp]
\centering\small
\begin{tabularx}{\linewidth}{@{}>{\raggedright\arraybackslash}p{0.17\linewidth}>{\raggedright\arraybackslash}X>{\raggedright\arraybackslash}X>{\raggedright\arraybackslash}p{0.21\linewidth}@{}}
\toprule
Item & Original reading & Correction & Source \\
\midrule
Jet oracle & ``simulable with $C=1$, closed at exponent $1/2$'', a derivation ``valid at all scales''; close jet candidates by theorem & 27 Jul (\texttt{CORR-20260727-003}): withdrawn, no derivation exists; closure countermanded; simulability open; \texttt{EV-GGM-002} supersedes \texttt{EV-GGM-001} & \art{crypto-autoresearcher:ledger/corrections/CORR-20260727-003.yaml} \\
\addlinespace
Planted generic-group null & Stage 1 null on planted $\ZZ/n$, $n$ of 16--20 bits & \texttt{EV-ECDLP-8a5106}: O-ARTIFACT, the kangaroo check M4 missed its $0.10$ band; decision \emph{refine} & \art{crypto-autoresearcher:ledger/evidence/EV-ECDLP-8a5106.yaml} \\
\addlinespace
ETH-route leads & twelve leads as written in \texttt{LEADS.md} & 2 Nov errata; told in \S\ref{sec:tn-corrections} & \art{research/ic_prime_eth_leads_20261031/LEADS.md} \S9 \\
\addlinespace
Interpolation of the discrete logarithm (Lange--Winterhof, Kiltz--Winterhof) & theorem statements read only in abstracts & 12 Nov: Kiltz--Winterhof 2004, Theorems 5 and 6 transcribed with their scope; Lange--Winterhof 2002 still unread (paywalled); \S\ref{sec:um-generic} & \art{research/cheon_aux_leads_20261112/LITERATURE_LANGE_WINTERHOF.md} \\
\addlinespace
Learned WDSat branching & learned order $16{,}097$ held-out conflicts against $3{,}092$ & confounded with the symmetry switch; told in \S\ref{sec:ln-components} & \art{cryptanalysis:experiments/ideas_learning_20261018/REPORT.md} \\
\addlinespace
Hashed-ECDL set-up hang & user report: Fourier-network construction impractical & 30 Oct: the AKS branch above \texttt{u64}; told in Table~\ref{tab:tn-native} & \art{research/hashed_ecdl_hash_prime_20261030/README.md} \\
\addlinespace
Eigen-S\'eta (not ECDLP) & \texttt{walk} variant; NIST categories 1/3/5 & 11 Oct: \texttt{walk} insecure (5/5 ciphertexts decrypted from public data); 17 Oct: categories hold only in gate count, unqualified claim withdrawn & \art{research/eigen_seta_prototype_20261011/README.md}, \art{research/eigen_seta_assessment_20261017/REPORT.md} \\
\bottomrule
\end{tabularx}
\caption[]{Corrections on unconventional routes, besides H-PSEUDO and xedni
(Table~\ref{tab:corr-xref}). All producer-reported.}
\label{tab:corr-exotic}
\end{table}

None of the rows in Table~\ref{tab:corr-exotic} opens or closes a route by
itself. Two changed a status in the other direction from the usual one: the
jet correction removed a closure that had no derivation behind it, and the
JKSST audit (Table~\ref{tab:corr-xref}) turned an exact-form closure into a
conditional one. A withdrawn closure is a correction, not evidence for the
route.

\section{Papers, notes and provenance}\label{sec:corr-papers}

Table~\ref{tab:corr-papers} lists corrections to papers and notes, and
provenance flags that bar a record from being cited as evidence.

\paragraph{The LLL degeneracy note.} The note reported two curve-specific
failure modes: secp256k1 failing LLL at 256 bits where P-256 and
brainpoolP256r1 succeed, and a uniform failure at 384 bits and above. Both
were overflow bugs in the \texttt{f64} Gram--Schmidt step ($n^4$ exceeds the
largest \texttt{f64}), and were withdrawn on 16 October
(\texttt{SHF-20261016-3caa6f}). On the fixed code every distinct seed pair at
256 and 384 bits recovers the key, secp256k1 included (3 of 3); P-521 at
$k_{\rm bits}=384$ recovers 0 of 3 on scaled \texttt{f64} and 3 of 3 with
2048-bit fixed-point Gram--Schmidt. On 17 October the mechanism sentence for
P-521 (``catastrophic cancellation'') was withdrawn as well
(\texttt{SHF-20261017-e4d7f6}), because with the scaling P-521's diagonal is
exactly $2^{500}$ and the described subtraction is exact
(\art{papers/lll-degeneracy-note/README.md}). The note has no index-calculus
content; it is listed because it is the clearest case here of a curve
``property'' that was a numerical bug.

\begin{table}[tbp]
\centering\small
\begin{tabularx}{\linewidth}{@{}>{\raggedright\arraybackslash}p{0.17\linewidth}>{\raggedright\arraybackslash}X>{\raggedright\arraybackslash}X>{\raggedright\arraybackslash}p{0.21\linewidth}@{}}
\toprule
Item & Original reading & Correction & Source \\
\midrule
LLL degeneracy note & secp256k1 and $\ge384$-bit curves fail LLL & withdrawn 16--17 Oct (see text) & \art{papers/lll-degeneracy-note/} \\
\addlinespace
Structural-completeness paper & ``the theorem is proven by case analysis'' & 17 Oct (\texttt{SHF-20261017-e4d7f6}, P-N4--N6): ``argued''; \S8 holds proof sketches; check counts recounted & \art{papers/structural-completeness/README.md} \\
\addlinespace
Isogeny-class rho paper & no test fired; inferential null & 19 Oct review: false (one B1 generator, two P1 cells, one P0 cell flagged; Kruskal--Wallis $p=0.039$); KS defect, missing tie corrections; null provisional, $p$-values not recomputed & \art{papers/isogeny-class-rho/REVIEW-20261019.md} \\
\addlinespace
Chung--Costello--Smith Recover & displayed Eqs.~(11), (12) & 22 Oct: in this implementation's convention the literal equations give 1{,}257 wrong signed divisors and 3 $D=0$ cases of 1{,}260; corrected equations recover all 1{,}260 & \art{research/kummer_recover_20261022/REPORT.md} \\
\addlinespace
\texttt{secpkk1} ladder (lit2026 W11(e)) & published 8/12/16-bit challenge cards & the cards use the opposite prime-search direction from the paper's own Algorithm~3 & \art{research/lit2026_w11_generic_engineering_20261011/README.md} \\
\addlinespace
Teske trapdoor evidence & \texttt{EV-ECTD-8c098d} rho-substitution prose & fabrication defect disclosed: the stated empirical justification has no supporting data & \art{crypto-autoresearcher:ledger/evidence/EV-ECTD-8c098d.yaml} \\
\addlinespace
Kernel-of-reduction design & \texttt{DEC-20261023-26a6de} cites \texttt{KN-LIT-7593} & the citation resolves to an unrelated AES paper & \art{crypto-autoresearcher:knowledge/leads/ecdlp/transfer/LEAD-TRF-kernel-of-reduction-lemma.yaml} \\
\addlinespace
Synthetic volcano study & volcano and ECDLP figures & invented data by its own banner; not evidence & \art{research/synthetic_volcano_ecdlp_20261018/README.md} \\
\addlinespace
Frobenius-quotient follow-up & reported SAT runs & supplied as a note; models and logs not committed; not reviewed evidence & \art{research/frobenius_quotient_followup_20260926/README.md} \\
\addlinespace
Scalar-3 Frobenius gate & new screen for $q=16$ & the mechanism is a rediscovery of the scalar-3/Jordan screen of \art{research/cubic_extension_transfer_20261025/}; only its extension to $q=16$ is new & \art{research/ghs_scalar3_frobenius_gate_20261027/REPORT.md} \\
\bottomrule
\end{tabularx}
\caption[]{Corrections to papers and notes, and provenance flags. All
producer-reported or sheriff-reported.}
\label{tab:corr-papers}
\end{table}

\paragraph{What the list shows.} Three patterns recur. First, a reference or a
pricing rule was corrected more often than an index-calculus count: the
class set-up, the strong Koblitz id, the negation map, the byte counter and
the transfer baseline each changed the distance to rho with no change to the
attacking side. Second, a preregistered control or gate that failed, as in the Cheon,
cofactor-clearing and E42 rows, voided the verdict or the ranking whatever the
measured quantity said; the rule was applied as frozen. Third, several
``properties'' of curves or classes (the LLL failures, the conductor-gap rho
shifts, the NIST coefficient test) turned out to be properties of a number
format, a coordinate model or a representation. As in
Chapter~\ref{ch:audit}, no correction listed here moved a generic or
leaderboard index-calculus row below a matched rho. The one row in this
appendix that ends below rho is the Joux--Vitse sieve route, which is restricted to the \texttt{jv\_cover} weak
class (full rational 2-torsion, never prime order) and is an accounting result
against a newly matched reference.
